\documentclass{article}

% lualatex handles UTF-8 natively; remove inputenc/fontenc for pdflatex compatibility
\usepackage{hyperref}
\usepackage{amsmath}
\usepackage{mathtools}
\usepackage{amssymb}
\usepackage{amsthm}
\usepackage{dsfont}
\usepackage{stmaryrd}
\usepackage{xcolor}
\usepackage{listings}
\usepackage{microtype}
\usepackage{fullpage}
\usepackage{authblk}
\usepackage{cleveref}

% ---- Theorem environments -----------------------------------------------
\newtheorem{theorem}{Theorem}[section]
\newtheorem{lemma}[theorem]{Lemma}
\newtheorem{definition}[theorem]{Definition}
\newtheorem{example}[theorem]{Example}
\newtheorem{remark}[theorem]{Remark}

% ---- Macros -------------------------------------------------------------
\newcommand{\R}{\mathbb{R}}
\newcommand{\Q}{\mathbb{Q}}
\newcommand{\N}{\mathbb{N}}
\newcommand{\Fin}[1]{\mathsf{Fin}(#1)}
\newcommand{\Flag}[1]{\mathcal{F}^{#1}}
\newcommand{\FlagAlg}[1]{\mathcal{A}^{#1}}
\newcommand{\poshom}[1]{\Phi^{#1}}
\newcommand{\den}[2]{p(#1,\,#2)}
\newcommand{\tdensity}[2]{\pi(#1;\,#2)}
\newcommand{\lean}[1]{\texttt{#1}}
\newcommand{\leanfmt}[1]{\texttt{\small #1}}

% Lean code style
\lstdefinelanguage{Lean4}{
  keywords={def,theorem,lemma,instance,structure,class,import,open,namespace,
             end,where,fun,let,have,show,by,match,with,if,then,else,
             return,do,for,in,noncomputable,abbrev,variable,section,
             native_decide,decide,norm_num,simp,ring,linarith,omega,
             exact,apply,rw,calc,intro,constructor,use,refine,obtain,
             rcases,cases,induction,assumption,contradiction,trivial,rfl},
  keywordstyle=\color{blue!70!black}\bfseries,
  comment=[l]{--},
  commentstyle=\color{gray}\itshape,
  stringstyle=\color{orange!80!black},
  morestring=[b]",
  sensitive=true,
  basicstyle=\ttfamily\small,
  breaklines=true,
  columns=fullflexible,
  escapeinside={(*@}{@*)},
  literate=
    {->}{{$\to$}}2
    {<-}{{$\leftarrow$}}2
    {<=}{{$\leq$}}1
    {>=}{{$\geq$}}1
    {alpha}{{$\alpha$}}1
    {sigma}{{$\sigma$}}1
    {forall}{{$\forall$}}1
    {exists}{{$\exists$}}1
    {phi}{{$\phi$}}1
    {psi}{{$\psi$}}1
}
\lstset{language=Lean4, frame=single, framesep=4pt,
        rulecolor=\color{black!20}, backgroundcolor=\color{gray!5}}

% =========================================================================
\title{Formalizing Flag Algebras in Lean~4 via Computational Reflection}
\date{}

\author[1]{Gyeongwon Jeong}
\affil[1]{School of Computing, KAIST, South Korea}

\author[2]{Seonghun Park}
\affil[2]{School of Computing, KAIST, South Korea}

\author[3]{Jihoon Hyun}
\affil[3]{School of Computing, KAIST, South Korea}

\author[4]{Sang-il Oum}
\affil[4]{Discrete Mathematics Group, Institute for Basic Science (IBS), South Korea}

\author[5]{Hongseok Yang}
\affil[5]{School of Computational Sciences, Korea Institute for Advanced Study (KIAS), South Korea}

% =========================================================================
\begin{document}
\maketitle

% =========================================================================
\begin{abstract}
Razborov's flag algebra method is one of the most powerful tools in extremal
combinatorics, having resolved many open problems about asymptotic subgraph
densities.  Applying it in practice, however, requires combining abstract algebraic
and measure-theoretic reasoning with large external computations: subgraph
density tables computed by enumeration, and semidefinite programming (SDP)
certificates found by numerical solvers.  Formalizing such proofs in a proof
assistant is therefore a challenge on two fronts: the abstract mathematical
structures must be faithfully encoded in a type theory, and the external
computational components must be imported and verified in a way that does not
compromise the overall proof's trustworthiness.

We present a Lean~4 formalization of Razborov's flag algebra method for
graphs, organized around a \emph{reflection-based} architecture.  At the
abstract level, we define flags as quotient types under graph isomorphism,
construct the flag algebra as a quotient module, equip it with a semantic
ordering via positive homomorphisms, and connect flag algebra inequalities to
combinatorial Turán densities through a measure-theoretic framework.  At the
computational level, we introduce a concrete, decidably-equal graph
representation (\lean{Sym2Graph}) and prove adequacy theorems connecting it to
the abstract definitions; this allows \lean{native\_decide} and \lean{decide
+kernel} to discharge hundreds of density and matrix-equality obligations
automatically.  Structural proof obligations---normalizing linear combinations
of flag terms and applying expansion and multiplication identities---are handled
by a suite of custom Lean~4 elaboration tactics that inspect the AST and exploit
a canonical naming convention for flags as a machine-readable encoding of
mathematical structure.

As results, we give complete formal proofs of Mantel's theorem and of the
Erd\H{o}s pentagon theorem ($\pi(K_3; C_5) = 24/625$), the latter including
both the upper bound via a formally verified SDP certificate and the lower
bound via an explicit blow-up construction.  To our knowledge, this is the
first formalization of the flag algebra method in any proof assistant.
\end{abstract}

% =========================================================================
\section{Introduction}
\label{sec:intro}

\paragraph{Flag algebras in extremal combinatorics.}
A central question in extremal graph theory is: among all large graphs on $n$
vertices that avoid some fixed graph $H$ as a (not necessarily induced) subgraph, what is the
maximum possible density of copies of another graph $F$?
The \emph{generalized Turán density} $\tdensity{F}{H}$
formalizes this question as a limit:
\[
  \tdensity{F}{H}
    \;=\;
  \lim_{n\to\infty}
  \frac{1}{\binom{n}{|V(F)|}}
  \max\bigl\{\#\text{copies of }F \text{ in } G
             \;\big|\;
             G \text{ is } H\text{-free},\;|V(G)|=n\bigr\}.
\]

Razborov's flag algebra method~\cite{razborov2007flag} provides a
\emph{systematic} framework for deriving upper bounds on such densities.
It works by constructing a graded, commutative $\mathbb{R}$-algebra---the
\emph{flag algebra}---whose elements represent linear combinations of
induced subgraph patterns.  The algebra is equipped with a product encoding
simultaneous occurrence and a semantic ordering: an element $f$ is
\emph{non-negative} (lies in the semantic cone) if its value under every
positive $\mathbb{R}$-algebra homomorphism is $\geq 0$.
Non-negativity certificates for carefully chosen elements yield upper bounds
on Turán densities directly.

In practice, these certificates are produced by semidefinite programming:
one seeks a positive semidefinite matrix $Q$ such that
$f - c\cdot\mathbf{1} = \sum_{i,j} Q_{ij}\cdot e_i e_j$ holds in the algebra
(where the $e_i$ are typed flag basis elements and $c$ is the claimed bound).
An SDP solver finds $Q$ numerically; one must then \emph{verify} that $Q$ is
genuinely positive semidefinite and that the algebraic identity holds exactly.

\paragraph{The formalization challenge.}
Formalizing a flag algebra proof in a proof assistant requires confronting three
distinct, non-reducible categories of proof obligation:

\begin{enumerate}
  \item \textbf{Abstract structure.}  Flags (quotients of labeled graphs under
    isomorphism), the flag algebra (a quotient module equipped with a commutative
    ring structure), positive homomorphisms (the semantic ordering),
    the Turán density (a measure-theoretic limit), and the transfer theorem
    connecting flag algebra inequalities to density bounds.  These are
    conceptually non-trivial but finite in number; each requires careful
    type-theoretic encoding.

  \item \textbf{Data-heavy computation.}  For each pair of flags $(F,G)$
    appearing in the proof, one must certify the exact rational density
    $\den{F}{G}$.  For the Erd\H{o}s pentagon theorem,
    this means hundreds of density values over graphs with up to five vertices
    and flag multiplication tables of similar size.  These values are computed
    externally and must be imported into the proof and verified against the
    formal definitions.

  \item \textbf{Algebraic bookkeeping.}  Flag algebra arguments involve
    manipulating linear combinations of hundreds of flag terms: expanding a
    flag at a larger vertex count, computing products of typed flags, and
    normalizing sums into a canonical form.  Each individual step is routine
    but the aggregate is prohibitively tedious to discharge manually.
\end{enumerate}

These three categories are not merely independent complications; they call for
\emph{qualitatively different} proof techniques.  The abstract structure
requires faithful encoding in a dependent type theory.
The data-heavy computation requires a \emph{reflection} architecture:
a decidably-computable concrete representation whose connection to the abstract
definitions is certified by adequacy theorems, allowing Lean's kernel (or native
evaluator) to check each density value automatically.
The algebraic bookkeeping requires \emph{proof-by-reflection via custom
elaboration tactics} that inspect the syntactic structure of the proof state
and dispatch the right sequence of domain-specific lemmas without manual
guidance.

\paragraph{This paper.}
We present a Lean~4 formalization of Razborov's flag algebra method for
graphs, organized around the two-layer architecture that the challenge analysis
dictates.  The \emph{abstract layer} encodes the mathematical semantics
faithfully in Lean~4's dependent type system, including a novel
\emph{general forbidden-subgraph reasoning rule} that does not require
problem-specific axiomatization.  The \emph{reflection layer} bridges the
abstract definitions to a finitely-computable concrete representation, enabling
automated discharge of density and SDP certificate obligations.

\paragraph{Contributions.}
\begin{itemize}
  \item \textbf{Abstract formalization (\S\ref{sec:abstract}).}
    We formalize the full abstract structure of Razborov's flag algebra in
    Lean~4: flags as quotient types of labeled graphs under isomorphism,
    the flag algebra as a quotient $\mathbb{R}$-module with a commutative ring
    structure, the semantic ordering via positive homomorphisms, and a
    measure-theoretic forbidden-subgraph reasoning framework
    (\lean{forbidLE}, \lean{generalizedTuranDensity\_le\_of\_forbidLE}).
    The forbidden-subgraph rule is formulated inside a \emph{single} ambient
    theory of simple graphs and does not require a per-problem axiom.

  \item \textbf{Reflection architecture (\S\ref{sec:reflection}).}
    We introduce \lean{Sym2Graph}, a finitely-representable graph type with
    decidable equality, and prove adequacy theorems equating abstract flag
    densities to computable \lean{Sym2Graph} densities.  We verify SDP
    certificates via an exact LDL$^\top$ decomposition over $\mathbb{Q}$,
    checked by \lean{decide +kernel}.  We distinguish a deliberate
    \emph{trust hierarchy}: \lean{native\_decide} for density tables (trusts
    the native compiler) and \lean{decide +kernel} for SDP certificates
    (trusts only the kernel).

  \item \textbf{Tactic automation (\S\ref{sec:tactics}).}
    We implement a suite of custom Lean~4 elaboration tactics that exploit a
    canonical naming convention for flag constants as a machine-readable
    encoding: \lean{ac\_sort\_pipeline} for linear normalization (replacing a
    generic sort that times out on expressions with $\sim\!25$ terms),
    and \lean{prove\_flag\_expand\_with\_forbidden\_flag} /
    \lean{prove\_flag\_mul\_with\_forbidden\_flag} for expansion and
    multiplication identities (each collapsing 15--20 manual steps into one
    tactic call).

  \item \textbf{Verified results (\S\ref{sec:results}).}
    We give formally complete proofs of Mantel's theorem and of
    $\tdensity{K_3}{C_5} = 24/625$ (the Erd\H{o}s pentagon theorem), including
    the upper bound via a formally verified SDP certificate and the lower bound
    via an explicit $C_5$-blow-up construction with a formal limit argument.
    To our knowledge, this is the first formalization of the flag algebra method
    in any proof assistant.
\end{itemize}

\paragraph{Paper organization.}
Section~\ref{sec:background} recalls the mathematical background on flag
algebras.  Sections~\ref{sec:abstract}--\ref{sec:tactics} describe the two
formalization layers in detail.  Section~\ref{sec:results} presents the
verified results.  Section~\ref{sec:related} discusses related work, and
Section~\ref{sec:conclusion} concludes.


\section{Background: Flag Algebras}
\label{sec:background}

This section recalls the mathematical definitions underlying flag algebras,
following Razborov~\cite{razborov2007flag}, and sets up the notation used
throughout the paper.
For $n \in \mathbb{N}$ we write $[n] = \{1,\ldots,n\}$.
All graphs are finite and simple.

\subsection{Types and Flags}

\paragraph{Types.}
A \emph{type} of size $k$ is a graph $\sigma$ with vertex set $[k]$.
The \emph{empty type} $\emptyset$ (size $k=0$) is the unique graph on the
empty vertex set; flags over $\emptyset$ are just ordinary finite graphs
up to isomorphism.

\paragraph{Flags.}
Fix a type $\sigma$ of size $k$.
A \emph{$\sigma$-flag} is a pair $G^\sigma = (G, \theta)$ where $G$ is a
finite graph and $\theta : [k] \hookrightarrow V(G)$ is an injective map such
that $\sigma$ is isomorphic to $G[\operatorname{Im}(\theta)]$ as a labeled
graph (i.e., $\theta$ is a graph embedding of $\sigma$ into $G$).
The integer $|V(G)|$ is the \emph{size} of $G^\sigma$.

Two $\sigma$-flags $(G_1,\theta_1)$ and $(G_2,\theta_2)$ are
\emph{isomorphic} if there is a graph isomorphism $\phi: V(G_1)\to V(G_2)$
with $\phi \circ \theta_1 = \theta_2$.  The set of isomorphism classes of
$\sigma$-flags of size $n$ is written $\mathcal{F}^\sigma_n$, and
$\mathcal{F}^\sigma = \bigcup_{n \geq k} \mathcal{F}^\sigma_n$.
Each $\mathcal{F}^\sigma_n$ is finite.

\subsection{Subflag Densities}

\paragraph{Single-flag density.}
For $\sigma$-flags $F$ of size $m$ and $G$ of size $n \geq m$, the
\emph{induced density} $\den{F}{G}$ is the probability that a uniformly
random injective map from $V(F) \setminus \operatorname{Im}(\theta_F)$ to
$V(G) \setminus \operatorname{Im}(\theta_G)$ (extending the type embedding)
induces a copy of $F$ in $G$ compatible with the type.  Explicitly:
\[
  \den{F}{G}
  \;=\;
  \frac{\bigl|\{\iota : V(F) \hookrightarrow V(G)
                  \mid \iota \text{ preserves type embedding and induces }F\}\bigr|}
       {\binom{n-k}{m-k}\,(m-k)!}.
\]
This value lies in $[0,1] \cap \mathbb{Q}$ and is invariant under
isomorphism of both $F$ and $G$.

\paragraph{Joint density.}
For $\sigma$-flags $F_1$ of size $m_1$ and $F_2$ of size $m_2$ and a host
flag $G$ of size $n \geq m_1 + m_2 - k$, the \emph{joint density}
$\den{F_1,F_2}{G}$ is the probability that two independently and uniformly
chosen injections from $V(F_i)\setminus [k]$ into $V(G)\setminus [k]$
(sampled without replacement from the same pool) each induce their
respective flag, compatible with the shared type embedding.

\subsection{The Flag Algebra}

Fix a type $\sigma$.  Let $\mathbb{R}[\mathcal{F}^\sigma]$ be the free
$\mathbb{R}$-module with basis $\mathcal{F}^\sigma$.
Define the \emph{zero space} $\mathcal{Z}^\sigma$ as the subspace generated
by all elements of the form
\[
  F - \sum_{G \in \mathcal{F}^\sigma_n} \den{F}{G} \cdot G,
  \qquad F \in \mathcal{F}^\sigma_m,\; n \geq m.
\]
The \emph{flag algebra} is the quotient module
\[
  \mathcal{A}^\sigma \;=\; \mathbb{R}[\mathcal{F}^\sigma] \,/\, \mathcal{Z}^\sigma.
\]
The zero-space relations express the combinatorial identity that, in a
sufficiently large graph, the average density of $F$ over all extensions of
the type embedding equals $\den{F}{G}$.

\paragraph{Multiplication.}
For $[F_1] \in \mathcal{A}^\sigma_{m_1}$ and $[F_2] \in \mathcal{A}^\sigma_{m_2}$,
their product at size $\ell \geq m_1 + m_2 - k$ is
\[
  [F_1] \cdot [F_2]
  \;=\;
  \sum_{G \in \mathcal{F}^\sigma_\ell} \den{F_1,F_2}{G} \cdot [G]
  \;\in\; \mathcal{A}^\sigma_\ell.
\]
This is well-defined on the quotient (independent of $\ell$) and makes
$\mathcal{A}^\sigma$ into a commutative, associative $\mathbb{R}$-algebra
with unit $[\text{type graph } \sigma]$.

\subsection{The Downward Operator and Semantic Non-Negativity}

\paragraph{Downward (unlabeling) operator.}
For a $\sigma$-flag $F$ of size $m$ (type of size $k$), the \emph{downward
operator} $\llbracket F \rrbracket_\sigma$ averages over all ways to embed
the type $\sigma$ into $F$, producing an element of the untyped algebra:
\[
  \llbracket F \rrbracket_\sigma
  \;=\;
  \frac{k!\,(m-k)!}{m!}
  \sum_{\theta: [k]\hookrightarrow V(F)}
  \bigl[(F, \theta)\bigr]_\emptyset,
\]
extended linearly to all of $\mathcal{A}^\sigma$.
The downward operator is an $\mathbb{R}$-module map
$\llbracket\cdot\rrbracket_\sigma : \mathcal{A}^\sigma \to \mathcal{A}^\emptyset$.

\paragraph{Positive homomorphisms.}
A \emph{positive homomorphism} is an $\mathbb{R}$-algebra homomorphism
$\phi: \mathcal{A}^\emptyset \to \mathbb{R}$ satisfying $\phi([G]) \geq 0$
for every graph $G$.
Positive homomorphisms correspond bijectively to convergent graph sequences:
every sequence $(G_n)$ with $|V(G_n)|\to\infty$ such that all flag densities
have limits determines a unique positive homomorphism~\cite{razborov2007flag}.

The \emph{semantic cone} is
$\mathcal{C} = \{f \in \mathcal{A}^\emptyset \mid \phi(f) \geq 0\ \forall \phi\}$.

\begin{theorem}[Razborov~{\cite{razborov2007flag}}]
  \label{thm:sdp-nonneg}
  For any type $\sigma$, flags $e_1,\ldots,e_r \in \mathcal{A}^\sigma$,
  and positive semidefinite matrix $A \in \mathbb{R}^{r\times r}$,
  \[
    \Bigl\llbracket \sum_{i,j=1}^r A_{ij}\, e_i\, e_j \Bigr\rrbracket_\sigma
    \;\in\; \mathcal{C}.
  \]
\end{theorem}

This theorem is the engine behind all SDP-based flag algebra bounds.
Given a target element $f_0 \in \mathcal{A}^\emptyset$ and a claimed bound
$c \in \mathbb{R}$, if one exhibits a PSD matrix $A$ and typed flags $e_i$ such
that
\[
  c \cdot \mathbf{1} - f_0
  \;=\;
  \Bigl\llbracket \sum_{i,j} A_{ij}\, e_i\, e_j \Bigr\rrbracket_\sigma,
\]
then $\phi(f_0) \leq c$ for all positive homomorphisms $\phi$, yielding an
upper bound on the Turán density of the corresponding graph pattern.

\subsection{Turán Densities and the Pentagon Problem}

\paragraph{Generalized Turán density.}
For unlabeled graphs $F$ and $H$, the \emph{generalized Turán density} is
\[
  \tdensity{F}{H}
  \;=\;
  \lim_{n\to\infty}
  \frac{\max\bigl\{|\text{copies of }F\text{ in }G|
                   \mid |V(G)|=n,\; G \text{ is }H\text{-free}\bigr\}}
       {\binom{n}{|V(F)|}}.
\]

\paragraph{The Erd\H{o}s pentagon problem.}
The problem asks for $\tdensity{C_5}{K_3}$: the maximum density of
pentagons in triangle-free graphs.  The answer $24/625$ was proved
independently by Grzesik~\cite{grzesik2012} and by Hatami, Hladk\'y, Kr\'al,
Norine, and Razborov~\cite{hatami2012}.  It is achieved by the blow-up of
$C_5$ (partition $n$ vertices into five nearly-equal groups and add all edges
between consecutive groups in the cycle).  We use this theorem as the main
case study throughout the paper, as it involves all three categories of proof
obligation: abstract structure, data-heavy density computation (hundreds of
rational values over 5-vertex graphs), and heavy algebraic bookkeeping.


\section{Abstract Formalization}
\label{sec:abstract}

The abstract formalization layer encodes the mathematical content of
Section~\ref{sec:background} faithfully in Lean~4's dependent type system.
The design principle is that every mathematical definition has a direct
Lean~4 counterpart that preserves the mathematical semantics exactly,
with no hidden invariants or definitional shortcuts.  This section describes
the four main components: flags as quotient types, the flag algebra as a
quotient module, the semantic ordering via positive homomorphisms, and a
general forbidden-subgraph reasoning framework.

\subsection{Flags as Quotient Types}

A \emph{type graph} of size $k$ is a \lean{SimpleGraph (Fin k)}, the standard
Lean~4/Mathlib type for simple graphs on a finite vertex set.  A
\emph{labeled graph} over type $\sigma$ with vertex set $V$ is a pair
\begin{lstlisting}
structure LabeledGraph (σ : FlagType (Fin n₀)) (V : Type) where
  graph     : SimpleGraph V
  typeEmbed : σ →gg graph    -- graph homomorphism embedding the type
\end{lstlisting}
where \lean{→gg} denotes a \lean{SimpleGraph} homomorphism that is also
injective on vertices.  Two labeled graphs over the same type are
\emph{flag-isomorphic} if they are related by a graph isomorphism that
commutes with the type embeddings.  The type of flags is then the quotient:
\begin{lstlisting}
def Flag (σ : FlagType (Fin n₀)) (V : Type) : Type :=
  Quotient (labeledGraphSetoid σ V)
\end{lstlisting}
The setoid \lean{labeledGraphSetoid} is defined from flag-isomorphism,
which is an equivalence relation on labeled graphs.

\paragraph{Size-indexed families.}
We write \lean{FlagWithSize σ n} for \lean{Flag σ (Fin n)} (flags of size
$n$) and \lean{FinFlag σ} for the dependent sum
$\Sigma\, (n : \mathbb{N}),\, \lean{FlagWithSize}\ \sigma\ n$, which collects
all flags across all sizes.
Each \lean{FlagWithSize σ n} is a \lean{Fintype} (inheriting finiteness from
the finiteness of \lean{Fin n} and the decidability of flag isomorphism,
established in Section~\ref{sec:reflection}), while \lean{FinFlag σ} is
\lean{Countable} and \lean{Infinite} when the type $\sigma$ is non-trivial.

\subsection{The Flag Algebra as a Quotient Module}

The flag algebra is built in
\lean{LeanFlagAlgebras.FlagAlgebra.FlagAlgebra}.
Following the mathematical definition, we first form the free
$\mathbb{R}$-module of finitely-supported functions on flags:
\begin{lstlisting}
abbrev FlagVector σ := FinFlag σ →₀ ℝ
\end{lstlisting}
(Here \lean{→₀} is Mathlib's \lean{Finsupp} type: functions with finite
support, forming a free module.)
The \emph{zero space} \lean{ZeroSpace σ} is defined as the submodule of
\lean{FlagVector σ} generated by all elements of the form
$[\lean{F}] - \sum_{G \in \lean{FlagWithSize}\ \sigma\ n} \den{F}{G} \cdot [G]$,
for each flag $F$ of size $m$ and each $n \geq m$.
The flag algebra is the quotient:
\begin{lstlisting}
def FlagAlgebra σ := FlagVector σ ⧸ ZeroSpace σ
\end{lstlisting}
We prove that \lean{FlagAlgebra σ} carries a commutative $\mathbb{R}$-algebra
structure, with product defined by
\begin{lstlisting}
def flagMulWithSize (F₁ : FlagWithSize σ m₁) (F₂ : FlagWithSize σ m₂)
    (ℓ : ℕ) (h : m₁ + m₂ - typeSize σ ≤ ℓ) : FlagVector σ :=
  ∑ G : FlagWithSize σ ℓ, flagDensity₂ F₁ F₂ G • unitVector G
\end{lstlisting}
and shown to be well-defined on the quotient via the zero-space relations.
The unit element is the class of the type graph $\sigma$ itself.

\subsection{Semantic Ordering via Positive Homomorphisms}

The ordering on \lean{FlagAlgebra σ} is semantic rather than syntactic.
We define:
\begin{lstlisting}
structure PositiveHom (σ : FlagType (Fin n₀)) where
  toFun       : FlagAlgebra σ →ₐ[ℝ] ℝ   -- R-algebra homomorphism
  nonneg      : ∀ F : FinFlag σ, 0 ≤ toFun ⟦unitVector F⟧
\end{lstlisting}
The \emph{semantic cone} and its induced ordering are:
\begin{lstlisting}
def semanticCone σ : Set (FlagAlgebra σ) :=
  {f | ∀ φ : PositiveHom σ, 0 ≤ φ.toFun f}

instance : Preorder (FlagAlgebra σ) :=
  { le := fun f g => g - f ∈ semanticCone σ, ... }
\end{lstlisting}

The main non-negativity theorem, proved in
\lean{LeanFlagAlgebras.FlagAlgebra.QuadraticForm}, is the Lean~4 counterpart
of Theorem~\ref{thm:sdp-nonneg}:
\begin{lstlisting}
theorem flagQuadraticForm_nonneg
    (M : Matrix (Fin r) (Fin r) ℝ) (hM : M.PosSemidef)
    (v : Fin r → FlagAlgebra σ) :
    0 ≤ ⟦ ∑ i j, M i j • (downward (v i * v j)) ⟧
\end{lstlisting}
This is the key lemma converting a PSD matrix into an element of the semantic
cone.  Its proof follows from the definition of positive homomorphisms and
the linearity of the downward operator: for any $\phi$,
$\phi\bigl(\llbracket \sum_{ij} M_{ij} v_i v_j \rrbracket\bigr)
= \mathbb{E}_\theta\bigl[\sum_{ij} M_{ij} \phi(v_i^{(\theta)}) \phi(v_j^{(\theta)})\bigr] \geq 0$
since $M$ is PSD and $\phi$ maps each basis vector to a non-negative real.

\subsection{The Forbidden-Subgraph Framework}
\label{sec:forbidden}

\paragraph{Motivation.}
In Razborov's original presentation~\cite{razborov2007flag}, applying the
flag algebra method to $H$-free graphs requires working in a modified theory
that axiomatizes the $H$-freeness condition from the outset.  This
per-problem axiomatization is pragmatically reasonable for a pen-and-paper
proof but is an obstacle to formalization: it would require a separate
verified theory for each forbidden subgraph.

\paragraph{Our approach.}
We instead develop a \emph{general forbidden-subgraph reasoning rule} that
works inside the ambient theory of simple graphs.  The key definition
(in \lean{LeanFlagAlgebras.Forbid.Basic}) is:
\begin{lstlisting}
def forbidLE (H : FinFlag ∅ₜ) (f g : FlagAlgebra σ) : Prop :=
  ∀ (φ₀ : PositiveHom ∅ₜ), φ₀.toFun ⟦unitVector H⟧ = 0 →
    ℙ[φ₀] {φ | φ.toFun f ≤ φ.toFun g} = 1
\end{lstlisting}
The notation $f \leq_{[H]} g$ means: almost surely under the measure
$\mathbb{P}^{\phi_0}$ (the ``random flag extension'' measure induced by any
positive homomorphism $\phi_0$ that assigns zero density to $H$), the typed
homomorphism satisfies $\phi(f) \leq \phi(g)$.

The measure $\mathbb{P}^{\phi_0}$ on typed positive homomorphisms is
constructed in \lean{LeanFlagAlgebras.FlagAlgebra.RandomHom}: the space of
positive homomorphisms for type $\sigma$ is given a compact topology (as a
closed subset of $\mathbb{R}^{\lean{FinFlag}\ \sigma}$, which is compact by
Tychonoff's theorem), and $\phi_0$ induces a Borel probability measure on it.

The transfer theorem is:
\begin{lstlisting}
theorem generalizedTuranDensity_le_of_forbidLE
    (H : FinFlag ∅ₜ) (F : FinFlag ∅ₜ) (c : ℝ) (hc : 0 < c)
    (hineq : forbidLE H (unitVector F) (c • 1)) :
    generalizedTuranDensity F H ≤ c
\end{lstlisting}
This single theorem, grounded in the measure-theoretic framework, is
completely general: any flag algebra inequality of the form
$f \leq_{[H]} c \cdot \mathbf{1}$ proved in the Lean theory immediately
yields a Turán density upper bound $\tdensity{F}{H} \leq c$, for any
forbidden subgraph $H$ and any target pattern $F$.  No additional
axiomatization is needed per application.

\paragraph{Design significance.}
This generality is a genuine formalization contribution, not a routine
consequence of the mathematical definitions.  The key technical work is
constructing the measure-theoretic framework (the random extension measure
$\mathbb{P}^{\phi_0}$), proving that it is well-defined as a probability
measure on the compact space of positive homomorphisms, and connecting it to
the combinatorial Turán density via a limit argument.  The result is
reusable: every future flag algebra proof in the library can use
\lean{generalizedTuranDensity\_le\_of\_forbidLE} as a black box.


\section{The Reflection Layer}
\label{sec:reflection}

The abstract formalization in \S\ref{sec:abstract} gives the right mathematical
definitions but is not directly evaluable: \lean{flagDensity} is defined via
non-constructive quotients, and \lean{Decidable} instances are not available
directly.  The reflection layer bridges this gap.

\subsection{Proof by Reflection in ITP}

The proof-by-reflection pattern in interactive theorem proving has three
components: (1) a \emph{concrete representation} over which a decision
procedure can evaluate; (2) an \emph{adequacy theorem} proving that the
decision procedure is sound with respect to the abstract definition; and (3)
a \emph{decision procedure} (the evaluator) that, given the concrete representation,
produces a Boolean answer which \lean{decide} converts into a proof.  Our
architecture applies this pattern at multiple granularities.

\subsection{Concrete Graph Representations}

We introduce \lean{Sym2Graph n} in
\lean{LeanFlagAlgebras.FlagAlgebra.Compute.Basic}:
\begin{lstlisting}
structure Sym2Graph (n : ℕ) where
  edges : Finset (Sym2 (Fin n))
  edges_valid : ∀ e ∈ edges, ¬e.IsDiag
\end{lstlisting}
This represents a graph on $\{0,\ldots,n-1\}$ as a finite set of unordered
non-diagonal pairs.  Unlike \lean{SimpleGraph (Fin n)} (which uses a
\lean{Prop}-valued adjacency relation), every operation on \lean{Sym2Graph n}
is computable: membership in \lean{edges} is decidable by \lean{Finset}
membership.

The same file derives the \lean{Fintype} instances needed to make graph
isomorphism decidable:
\begin{lstlisting}
instance : Fintype (G₁ ↪g G₂) := ...   -- graph embeddings
instance : Fintype (G₁ ≃g G₂) := ...   -- graph isomorphisms
instance : Decidable (Nonempty (G ≃f G')) := ...  -- flag isomorphism
\end{lstlisting}
These instances are engineered carefully: they enumerate candidates (all
injections from $V(G_1)$ to $V(G_2)$, or all bijections) and filter those
that preserve adjacency and type embeddings.  This reduces the question ``are
these two labeled graphs isomorphic?'' to a finite, decidable search.

\subsection{Adequacy Theorems}

The connection between the abstract and concrete worlds is established by
adequacy theorems in \lean{LeanFlagAlgebras.FlagAlgebra.Compute.FlagDensity}:

\begin{lstlisting}
theorem flagDensity₁_eq_sym2EmptyTypeFlagDensity₁
    (F : Sym2EmptyTypedFlag m) (G : Sym2EmptyTypedFlag n) :
    flagDensity₁ F.toFlag G.toFlag =
    sym2EmptyTypeFlagDensity₁ F G

theorem flagDensity₂_eq_sym2EmptyTypeFlagDensity₂
    (F₀ : Sym2EmptyTypedFlag m₀) (F₁ : Sym2EmptyTypedFlag m₁)
    (G : Sym2EmptyTypedFlag n) :
    flagDensity₂ F₀.toFlag F₁.toFlag G.toFlag =
    sym2EmptyTypeFlagDensity₂ F₀ F₁ G
\end{lstlisting}
The right-hand sides---\lean{sym2EmptyTypeFlagDensity₁} and
\lean{sym2EmptyTypeFlagDensity₂}---are defined entirely over \lean{Sym2Graph}
and are fully computable.  After rewriting a density goal with these theorems,
the goal becomes a decidable proposition about a finite computation over
\lean{Finset}, which \lean{native\_decide} can evaluate directly.

\subsection{Elaboration-Time Theorem Generation}

Having established the adequacy theorems, we need to generate and prove the
density and multiplication lemmas needed for each specific theorem.  Rather
than writing these by hand, we use Lean~4's elaboration monad to generate them
automatically from precomputed data.

The elaboration macro \lean{load\_flag\_pair\_density\_theorems} in
\lean{LeanFlagAlgebras.ErdosPentagon.Densities.DensityLoader} reads a JSON
file containing density values, and for each entry \lean{(F1, F2, G, p/q)}
generates and immediately proves the theorem:
\begin{lstlisting}
@[simp]
theorem flagDensity₂_Flag_5_1_0_3_Flag_5_1_0_7_Flag_5_0_0_2
    : flagDensity₂ Flag_5_1_0_3 Flag_5_1_0_7 Flag_5_0_0_2 = 3/10
  := by
  delta Flag_5_1_0_3 Flag_5_1_0_7 Flag_5_0_0_2
  rw [flagDensity₂_eq_sym2FlagDensity₂]
  native_decide
\end{lstlisting}
The three proof steps are: unfold the concrete flag definitions (making their
\lean{Sym2Graph} structure explicit), rewrite with the adequacy theorem, then
call \lean{native\_decide} to evaluate the resulting finite computation.
Hundreds of such theorems are generated and proved at compile time.  The same
elaboration pattern is used by \lean{MulLoader} for flag multiplication tables.

\subsection{SDP Certificate Verification via LDL$^\top$}
\label{sec:sdp}

The upper bound proof for the Erd\H{o}s pentagon theorem requires showing
that a sum of three quadratic forms (one for each 1-vertex type) is
non-negative under the $C_5$-free assumption.  The certificates are three
positive semidefinite matrices $P$, $Q$, $R$ over $\Q$ (each $8 \times 8$),
found externally by an SDP solver.

We verify positive semidefiniteness via an LDL$^\top$ decomposition.  For
matrix $P$, we exhibit a lower-triangular matrix $L_P$ and a diagonal matrix
$D_P$ with non-negative diagonal entries $d_P$, defined as explicit rational
constants in \lean{LeanFlagAlgebras.ErdosPentagon.Matrix.PosSemiDef}:
\begin{lstlisting}
def P : Matrix (Fin 8) (Fin 8) ℚ := !![(24/625 : ℚ), (-36/625 : ℚ), ...]
def dP : Fin 8 → ℚ := ![(24/625 : ℚ), (223/625 : ℚ), ..., 0, 0, 0, 0]
def LP : Matrix (Fin 8) (Fin 8) ℚ := !![(1 : ℚ), 0, ...; (-3/2 : ℚ), 1, ...]
\end{lstlisting}
Positive semidefiniteness then follows from two lemmas:
\begin{lstlisting}
lemma dP_nonneg (i : Fin 8) : 0 ≤ dP i := by
  fin_cases i <;> norm_num [dP]

lemma P_eq_LDL : P = LP * Matrix.diagonal dP * LPᵀ := by
  decide +kernel

theorem P_posSemidef : P.PosSemidef :=
  posSemidef_of_eq_mul_diagonal_mul_transpose dP_nonneg P_eq_LDL
\end{lstlisting}
The matrix equality \lean{P\_eq\_LDL} is a decidable proposition over $\Q$
(rational matrix equality), verified by \lean{decide +kernel}---the Lean
kernel evaluates the matrix product symbolically and checks equality entry by
entry.  This uses no compiled native code, so no axiom beyond the standard
Lean kernel is assumed.

\paragraph{The trust hierarchy.}
We deliberately use two different decision procedures with different trust
profiles:
\begin{itemize}
  \item \lean{native\_decide} compiles the goal to native code and evaluates
    it.  It relies on the correctness of the Lean-to-native compiler, which is
    outside the kernel.  We use it for density tables, where the computation is
    large but the individual claims are simple (a rational equality between the
    defined density and a stated value).

  \item \lean{decide +kernel} evaluates inside the kernel.  It is slower but
    uses no trusted computation outside the Lean kernel itself.  We use it for
    the SDP matrix equalities, where trust is most critical.
\end{itemize}
In both cases, the external computation (the density enumeration scripts, the
SDP solver) is responsible only for producing \emph{candidates}---values to
claim.  The Lean proof is responsible for verifying each claim against the
formal definition.

% =========================================================================
\section{The Tactic Layer}
\label{sec:tactics}

The reflection layer handles data-heavy computation.  A complementary layer of
custom elaboration tactics handles the structurally repetitive proof steps that
arise in every flag algebra argument.

\subsection{Naming Convention as Machine-Readable Encoding}

A key enabling design decision is that every flag and flag algebra constant is
named according to the scheme \lean{Flag\_n\_k\_m\_i} and
\lean{FlagAlgebra\_n\_k\_m\_i}, where $n$ is the vertex count, $k$ and $m$
identify the type graph, and $i$ is the canonical index within that
isomorphism class.  This convention is not merely organizational: it serves as
a machine-readable encoding of mathematical structure that all custom tactics
exploit.  By parsing the trailing integers from a constant's name, a tactic can
determine:
\begin{itemize}
  \item the \emph{canonical order} of a flag term in a linear combination,
  \item the \emph{precomputed lemma name} for its density or multiplication
    identity, and
  \item the \emph{flag type parameters} needed to instantiate the appropriate
    adequacy theorem.
\end{itemize}
This design is similar to the use of canonical forms in certified compilation:
the naming scheme is a ``normal form'' that tactics can compute with directly,
without needing to reduce the mathematical content of the expression.

\subsection{Linear Normalization: \lean{ac\_sort\_pipeline}}

\paragraph{The problem.}
After expanding a flag algebra element at a larger size, the proof obligation
is typically an equality between two linear combinations of flag terms that are
provably equal but syntactically in different orders.  Generic tactics such as
\lean{ring} or \lean{ac\_rfl} handle this in principle but become prohibitively
slow on large expressions ($\geq 20$ terms) because they must search over all
permutations of summands.

\paragraph{The solution.}
The \lean{sort} family of tactics, implemented in
\lean{LeanFlagAlgebras.ErdosPentagon.SortTactic}, normalizes a linear
combination by:
\begin{enumerate}
  \item \emph{Flattening.}  The tactic traverses the expression's AST,
    pattern-matching on \lean{HAdd.hAdd}, \lean{HSub.hSub},
    \lean{HSMul.hSMul}, and \lean{Neg.neg}, to produce a flat list of
    \lean{(base, coefficient)} pairs.
  \item \emph{Sorting.}  Each base term is assigned a sort key: the trailing
    integer in its constant name (e.g., \lean{FlagAlgebra\_5\_1\_0\_3} gets
    key $3$), with the pretty-printed string as a tiebreaker.  The list is
    sorted by this key using insertion sort.
  \item \emph{Rebuilding.}  The sorted list is reassembled into a Lean
    expression (left-associative sum of scalar multiples).
  \item \emph{Equality proof.}  The tactic proves that the original and sorted
    expressions are equal by \lean{ac\_rfl} or \lean{abel\_nf}, then replaces
    the goal with the sorted form.
\end{enumerate}
The full pipeline \lean{ac\_sort\_pipeline} additionally runs \lean{norm\_num}
before and after sorting, and collects like terms with \lean{$\leftarrow$ add\_smul}
after sorting, so that the resulting expression is fully simplified.

\paragraph{Why this is necessary, not cosmetic.}
The comment in \lean{ErdosPentagon.Lemmas} is revealing:
``\texttt{-- sort\_at -- takes more than 10 minutes}''.
The generic \lean{sort\_at} (which uses \lean{simp} without index guidance)
times out on expressions with $\sim 25$ flag terms.  The index-guided
\lean{ac\_sort} completes the same step in under a second, by exploiting domain
knowledge to bypass combinatorial search.

\subsection{Flag Expansion and Multiplication Tactics}

\paragraph{The problem.}
A central step in every flag algebra proof is the \emph{expansion} identity: a
small flag $F \in \Flag{\sigma}_m$, viewed in the algebra at size $n > m$,
equals a weighted sum of all $n$-vertex flags.  Similarly, the product of two
typed flags equals a linear combination at the product size.  These identities
must be verified for each specific flag and expansion size that appears in the
proof---dozens of instances per theorem.

\paragraph{\lean{prove\_flag\_expand\_with\_forbidden\_flag} N.}
This tactic, implemented in \lean{LeanFlagAlgebras.Logic.Tactic}, proves goals
of the form:
\[
  \bigl(\lean{FlagAlgebra\_3\_0\_0\_3} =_{a} 0\bigr)
  \;\vdash_a\;
  \lean{FlagAlgebra\_2\_0\_0\_1} =_{a} \tfrac{1}{3} \cdot \lean{FlagAlgebra\_3\_0\_0\_1}
    + \tfrac{2}{3} \cdot \lean{FlagAlgebra\_3\_0\_0\_2}
\]
where $=_{a}$ is the assertion operator in the logical framework.  The tactic:
\begin{enumerate}
  \item Scans the goal's AST for all \lean{FlagAlgebra\_*} and \lean{Flag\_*}
    constants using \lean{collectPrefixConstants}, and parses their indices
    $(n, k, m, i)$ using \lean{parseFlagAlgebraIndices?}.
  \item Unfolds the flag definitions by name, introduces the positive
    homomorphism $\phi$ and the forbidden-flag hypothesis $h$.
  \item Applies \lean{unitVector\_quot\_eq\_sum} to expand the flag at size
    $N$, then rewrites with the preloaded lemmas \lean{flagSet\_N\_k\_m\_eq\_univ}
    and \lean{flagSet\_N\_k\_m\_val\_eq} (themselves generated by the loaders).
  \item Simplifies the expansion using $h$ (the forbidden flag has value~0)
    and closes with \lean{ring\_nf}.
\end{enumerate}

\paragraph{\lean{prove\_flag\_mul\_with\_forbidden\_flag} N.}
This tactic handles product identities under a forbidden-flag assumption.  It
follows the same structure but applies \lean{unitVector\_quot\_mul\_eq\_flagMul\_quot}
and unfolds the flag multiplication at size $N$.  It also normalizes the
multiplication order: if the two operands' indices are in non-canonical order,
it first applies \lean{mul\_comm} before the rest of the proof, ensuring that
precomputed lemmas (which assume a fixed canonical ordering) can be applied
directly.

The result of these two tactics is that expansion and multiplication identities
that would individually require 15--20 tactic steps are each proved by a single
tactic invocation:
\begin{lstlisting}
example : FlagAlgebra_3_0_0_3 =ₐ 0
    ⊢ₐ FlagAlgebra_2_0_0_1 =ₐ (1/3 : ℝ) • FlagAlgebra_3_0_0_1
                             + (2/3 : ℝ) • FlagAlgebra_3_0_0_2
  := by prove_flag_expand_with_forbidden_flag 3
\end{lstlisting}

% =========================================================================
\section{Results}
\label{sec:results}

Using the formalization infrastructure described in the preceding sections,
we give formally complete proofs of two theorems in extremal combinatorics.
Both proofs are machine-checked in Lean~4 and contain no \lean{sorry}
placeholders in their proof paths.

\subsection{Mantel's Theorem}

Mantel's theorem states that a triangle-free graph on $n$ vertices has at
most $\lfloor n^2/4 \rfloor$ edges, achieved by the complete bipartite graph
$K_{\lfloor n/2\rfloor, \lceil n/2\rceil}$.  In flag algebra terms:

\begin{lstlisting}
theorem Mantel_theorem : K2 ≤ (1/2 : ℝ) • 1 + K3
\end{lstlisting}

Here $K_2$ and $K_3$ are elements of \lean{FlagAlgebra ∅ₜ} representing the
single edge and the triangle.  The inequality asserts that in every sequence
of graphs whose triangle density tends to zero, the edge density is at most
$1/2$.

The proof uses an explicit square in the algebra of 1-vertex-typed flags.
Let $a = \lean{FlagAlgebra\_2\_1\_0\_0}$ and $b = \lean{FlagAlgebra\_2\_1\_0\_1}$
be the two typed flags of size 2 with a 1-vertex type (corresponding to
the two labeled graphs on two vertices with one labeled vertex: one where the
two vertices are adjacent, one where they are not).
Then
\[
  \bigl\llbracket (a - b)^2 \bigr\rrbracket_1
  \;=\;
  2 \cdot K_2 - \mathbf{1}
  \;\geq\; 0
  \;\;\text{in the semantic cone.}
\]
This is enough to conclude $K_2 \leq \frac{1}{2} \cdot \mathbf{1}$ in the
presence of $K_3 = 0$.  The SDP certificate here is small enough to write
explicitly (no external solver needed), making Mantel's theorem a clean
end-to-end illustration of the method before the heavier machinery of the
pentagon proof.

The formal proof in \lean{LeanFlagAlgebras.MantelTheorem.MantelTheorem} uses
\lean{prove\_flag\_expand\_with\_forbidden\_flag} to establish the expansion
identity, \lean{ac\_sort\_pipeline} to normalize the resulting linear
combination, and \lean{flagQuadraticForm\_nonneg} to confirm the semantic
non-negativity.

\subsection{The Erd\H{o}s Pentagon Theorem}

The main result is:
\begin{lstlisting}
theorem ErdosPentagon_Turan
    : generalizedTuranDensity C5 K3 = 24/625
\end{lstlisting}
proved as the conjunction of an upper bound and a lower bound.

\paragraph{Upper bound.}
\begin{lstlisting}
theorem ErdosPentagon_Turan_upperBound
    : generalizedTuranDensity C5 K3 ≤ 24/625 :=
  generalizedTuranDensity_le_of_forbidLE (by norm_num)
    ErdosPentagon_flagAlgebra
\end{lstlisting}
The key lemma \lean{ErdosPentagon\_flagAlgebra} establishes the flag algebra
inequality $C_5 \leq_{[K_3]} \frac{24}{625} \cdot \mathbf{1}$, i.e., that
$C_5$-density is at most $24/625$ under the $K_3$-free assumption.
The proof exhibits three PSD matrices $P, Q, R$ over $\mathbb{Q}$
(one per 1-vertex type, each $8\times 8$) and shows that the sum of the
corresponding quadratic forms, after downward-averaging, equals
$\frac{24}{625} \cdot \mathbf{1} - [C_5]$.

Positive semidefiniteness of each matrix is verified via the LDL$^\top$
approach described in Section~\ref{sec:sdp}.  Once PSD is established, the
quadratic form non-negativity follows from
\lean{flagQuadraticForm\_nonneg}.  The equality between the quadratic form
sum and the claimed bound is verified by the tactic pipeline:
\lean{prove\_flag\_expand\_with\_forbidden\_flag} and
\lean{prove\_flag\_mul\_with\_forbidden\_flag} handle the expansion and
product identities, while \lean{ac\_sort\_pipeline} normalizes the resulting
expressions.

The final one-line proof \lean{ErdosPentagon\_Turan\_upperBound} connects
\lean{ErdosPentagon\_flagAlgebra} to the combinatorial statement via
\lean{generalizedTuranDensity\_le\_of\_forbidLE}.

\paragraph{Lower bound.}
The lower bound
\begin{lstlisting}
theorem ErdosPentagon_Turan_lowerBound
    : generalizedTuranDensity C5 K3 ≥ 24/625
\end{lstlisting}
is proved by an explicit construction.  We define the blow-up of a graph:
\begin{lstlisting}
def blowUp (G : SimpleGraph V) (n : ℕ) : SimpleGraph (V × Fin n) :=
  { Adj := fun v w => G.Adj v.1 w.1 ∧ v ≠ w, ... }
\end{lstlisting}
The $n$-fold blow-up of $C_5$ (replacing each vertex by an independent set
of $n$ vertices and each edge by a complete bipartite graph) is triangle-free:
\begin{lstlisting}
theorem blowUp_K3_free : (blowUp C5 n).CliqueFree 3
\end{lstlisting}
It has $5n$ vertices and at least $n^5$ induced copies of $C_5$:
\begin{lstlisting}
theorem subgraphCount_blowUp_C5_ge (n : ℕ) :
    n^5 ≤ (blowUp C5 n).inducedSubgraphCount C5
\end{lstlisting}
The lower bound on the Turán density follows from the limit:
\[
  \lim_{n\to\infty} \frac{n^5}{\binom{5n}{5}} = \frac{24}{625},
\]
which is proved as an exact algebraic identity using \lean{norm\_num} after
unfolding the definition of $\binom{5n}{5} = \frac{5n(5n-1)(5n-2)(5n-3)(5n-4)}{120}$.
The limit argument is formalized using Lean's \lean{Filter.Tendsto} framework:
\begin{lstlisting}
lemma tendsto_C5_blowUp_density :
    Filter.Tendsto (fun n => n^5 / Nat.choose (5*n) 5)
                   Filter.atTop (nhds (24/625))
\end{lstlisting}

\paragraph{Trust hierarchy.}
The proof of \lean{ErdosPentagon\_Turan} depends on two classes of
computational verification:
\begin{itemize}
  \item \lean{native\_decide} is used for all density table equalities
    (e.g., $\den{F_1,F_2}{G} = p/q$ for specific flags).  This tactic
    compiles the goal to native code and evaluates it; it relies on the
    correctness of the Lean-to-native compiler, which is outside the kernel.
    We accept this trust assumption because each individual density claim is
    a simple rational equality, and the native evaluator has been extensively
    tested in the Lean community.

  \item \lean{decide +kernel} is used for the three SDP matrix equalities
    $P = L_P D_P L_P^\top$, $Q = L_Q D_Q L_Q^\top$, $R = L_R D_R L_R^\top$.
    This tactic evaluates inside the Lean kernel using no compiled native
    code, so it introduces no trust beyond the standard Lean axioms
    (\lean{propext}, \lean{Quot.sound}, \lean{Classical.choice}).
    We use \lean{decide +kernel} specifically here because the SDP
    certificates are the trust-critical component: if a matrix is claimed PSD
    but is not, the entire upper bound proof collapses.
\end{itemize}
In both cases, the external computation (density enumeration scripts
in \lean{LeanFlagAlgebras/ErdosPentagon/Densities/} and the SDP solver in
\lean{LeanFlagAlgebras/ErdosPentagon/Matrix/}) is responsible only for
producing \emph{candidate} values.  The Lean proofs are responsible for
verifying each candidate against the formal definition.  There are no
\lean{sorry}s or \lean{axiom}s in the proof paths of either
\lean{Mantel\_theorem} or \lean{ErdosPentagon\_Turan}.


\section{Related Work}
\label{sec:related}

\paragraph{Formalizations of combinatorics.}
Several combinatorial theorems have been formalized in Lean and related
systems.  Dillies and Mehta~\cite{dillies2022szemeredi} formalized
Szemerédi's Regularity Lemma in Lean~4; Mehta~\cite{mehta2022kruskal}
formalized Kruskal-Katona; Subercaseaux et al.~\cite{subercaseaux2024hexagon}
formally verified the empty hexagon number using a SAT-based pipeline.  Our
work differs in that the central method (flag algebras) relies on an algebraic
quotient construction and measure-theoretic limits, rather than purely
combinatorial enumeration.

\paragraph{Computer-assisted flag algebra arguments.}
The Flagmatic software system~\cite{razborov2010flagmatic} automates the
computation of flag algebra certificates, but produces results that are
validated only informally.  Our work provides the first formally verified
end-to-end pipeline from flag algebra certificates to combinatorial theorems.

\paragraph{Proof by reflection.}
Proof by reflection is a classical technique in the Coq community (see
Chlipala~\cite{chlipala2013cpdt} for a survey), and has been used in
Lean for verified computation (e.g., in Mathlib's \lean{decide} and
\lean{norm\_num} tactics).  Our use of \lean{Sym2Graph} as a concrete mirror
of \lean{SimpleGraph} follows the same pattern, extended to the domain of flag
density computation.

\paragraph{Tactic meta-programming.}
Custom tactics for domain-specific proof automation have been developed in many
contexts~\cite{ebner2017structured,sozeau2008coq}.  Our tactics exploit the
naming convention as a form of \emph{reflected metadata}: the name of a
constant encodes information that the tactic exploits without having to reduce
the mathematical content of the expression.  This is related to, but distinct
from, approaches based on type-class inference or user-facing attribute
annotations.

% =========================================================================
\section{Conclusion}
\label{sec:conclusion}

We have presented the first formalization of Razborov's flag algebra method in
a proof assistant, realized in Lean~4.  The formalization is organized around
a two-layer architecture.  The \emph{reflection layer} connects the abstract
algebraic definitions to a concrete, decidably-computable graph representation
via adequacy theorems, enabling \lean{native\_decide} and \lean{decide +kernel}
to automatically discharge hundreds of density and SDP certificate obligations.
The \emph{tactic layer} provides custom elaboration tactics---exploiting a
canonical naming scheme for flag constants---that handle the structural
bookkeeping (linear normalization, expansion identities, multiplication
identities) of flag algebra proofs.

The end results are formally complete proofs of Mantel's theorem and the
Erd\H{o}s pentagon theorem ($\tdensity{K_3}{C_5} = 24/625$).  We believe
the architecture generalizes: the $\sim 100$ known flag algebra results in
extremal combinatorics all involve the same three categories of proof
obligation (abstract structure, data-heavy computation, and algebraic
bookkeeping), and our framework provides reusable infrastructure for all three.

\paragraph{Future work.}
Replacing \lean{native\_decide} with \lean{decide +kernel} throughout (or
with a formally verified external checker) would eliminate the remaining
dependency on the native compiler.  Extending the framework beyond graphs---to
hypergraphs, directed graphs, or other combinatorial structures---would require
generalizing the type-parameter conventions but no new conceptual machinery.
Automating the discovery of SDP certificates from within Lean (rather than
importing them from external solvers) remains a longer-term goal.

% =========================================================================
\bibliographystyle{plain}
\begin{thebibliography}{99}

\bibitem{razborov2007flag}
A.~A. Razborov.
\newblock Flag algebras.
\newblock \textit{Journal of Symbolic Logic}, 72(4):1239--1282, 2007.

\bibitem{grzesik2012}
A.~Grzesik.
\newblock On the maximum number of five-cycles in a triangle-free graph.
\newblock \textit{Journal of Combinatorial Theory, Series B},
  102(5):1061--1066, 2012.

\bibitem{hatami2012}
H.~Hatami, J.~Hladk\'{y}, D.~Kr\'{a}l, S.~Norine, and A.~Razborov.
\newblock On the number of pentagons in triangle-free graphs.
\newblock \textit{Journal of Combinatorial Theory, Series A},
  120(3):722--732, 2013.

\bibitem{dillies2022szemeredi}
Y.~Dillies and B.~Mehta.
\newblock Formalising Szemer\'{e}di's regularity lemma in Lean.
\newblock In \textit{Proc.\ ITP 2022}, LIPIcs~237, 2022.

\bibitem{mehta2022kruskal}
B.~Mehta.
\newblock Formalising the Kruskal-Katona theorem in Lean.
\newblock In \textit{Proc.\ ITP 2022}, 2022.

\bibitem{subercaseaux2024hexagon}
B.~Subercaseaux, M.~J.~H. Heule, J.~Mackey, J.~Meadows, R.~Tao, and
  C.~Wu.
\newblock Formal verification of the empty hexagon number.
\newblock In \textit{Proc.\ ITP 2024}, LIPIcs~309, 2024.

\bibitem{razborov2010flagmatic}
A.~A. Razborov.
\newblock On 3-hypergraphs with forbidden 4-vertex configurations.
\newblock \textit{SIAM Journal on Discrete Mathematics}, 24(3):946--963, 2010.

\bibitem{chlipala2013cpdt}
A.~Chlipala.
\newblock \textit{Certified Programming with Dependent Types}.
\newblock MIT Press, 2013.

\bibitem{ebner2017structured}
G.~Ebner, S.~Ullrich, J.~Roesch, J.~Avigad, and L.~de Moura.
\newblock A metaprogramming framework for formal verification.
\newblock \textit{Proc.\ ACM Program.\ Lang.}, 1(ICFP):34:1--34:29, 2017.

\bibitem{sozeau2008coq}
M.~Sozeau and N.~Oury.
\newblock First-class type classes.
\newblock In \textit{Proc.\ TPHOLs 2008}, LNCS~5170, 2008.

\end{thebibliography}

\end{document}
